\documentclass[a4paper]{article}

\usepackage{graphicx}
\usepackage{amsmath,amssymb}
\usepackage{minted}
\usepackage{multirow}
\usepackage{float}
\usepackage{todonotes}
\usepackage{forest}
\usepackage{xstring}
\usepackage{xcolor}
\usepackage[most]{tcolorbox}
\usepackage{xparse}
\usepackage{natbib}

\usetikzlibrary{graphs,graphdrawing}
\usegdlibrary{layered}

% for external graphviz
\input{/home/donkarlo/Dropbox/repo/nd_language_ai_project/src/nd_language_ai/formal/computer/programming/tex/latex/code_clips/external_graphviz.tex}

% for tags
\input{/home/donkarlo/Dropbox/repo/nd_language_ai_project/src/nd_language_ai/formal/computer/programming/tex/latex/part/control_sequence/kind/semtag/semtag.tex}

% for links
\input{/home/donkarlo/Dropbox/repo/nd_language_ai_project/src/nd_language_ai/formal/computer/programming/tex/latex/code_clips/seehere.tex}

\title{Memory trace or engram}
\date{}
\author{Mohammad Rahmani}

\begin{document}
   \maketitle

   \section{Definition}
      \section{Definition}
      An engram, or memory trace, is an enduring, experience-induced physical and/or chemical modification distributed across a specific neuronal ensemble that constitutes the neural representation of a memory.
      The ensemble consists of neurons recruited and modified during learning, and subsequent reactivation of these neurons contributes to, or can trigger, retrieval of the corresponding memory \cite{josselyn-2020-memory-engrams-recalling-the-past-and-imagining-the-future,sakaguchi-2012-catching-the-engram-strategies-to-examine-the-memory-trace}.

      An engram is a unit of cognitive information imprinted in a physical substance, theorized to be the means by which memories are stored as biophysical or biochemical changes in the brain or other biological tissue, in response to external stimuli \seeherefooter{https://en.wikipedia.org/wiki/Engram_(neuropsychology)}.

      \section{Paperas about engram}
      \cite{josselyn-2020-memory-engrams-recalling-the-past-and-imagining-the-future} reviews the modern concept of the engram as an enduring neural representation produced by an experience. It presents observational, loss-of-function, gain-of-function, and mimicry experiments that provide causal evidence for the existence of engram cells. Neuronal excitability influences which neurons are allocated to an engram during memory formation. Persistent synaptic plasticity, altered connectivity, and different active or silent states contribute to the lifetime and accessibility of an engram. The authors ultimately propose the engram, including distributed brain-wide engram complexes, as a fundamental unit underlying memory storage and retrieval.


      \cite{sakaguchi-2012-catching-the-engram-strategies-to-examine-the-memory-trace} reviews experimental strategies for identifying the physical neuronal traces that represent memories in the brain. It describes how activity-dependent immediate early genes such as \textit{c-fos} and \textit{Arc} can be used to identify neuronal ensembles recruited during learning. Genetic techniques make it possible to permanently label neurons involved in a particular experience and examine their later reactivation.
      Selective manipulation of these neuronal populations provides evidence that particular ensembles participate causally in memory encoding and retrieval. The review therefore connects the classical concept of the memory trace with experimentally identifiable and manipulable neuronal ensembles at cellular and circuit levels.
      
      
      

      \bibliographystyle{plainnat}
      \bibliography{/home/donkarlo/Dropbox/repo/nd_shared_research/bibliography/bibliography.bib}
\end{document}